\title{FlashAttention-3: Fast and Accurate Attention with Asynchrony and Low-precision}

\begin{abstract}
As a central component within the widespread Transformer architecture, attention mechanisms remain a primary computational constraint in large language models and applications demanding extended context lengths. \textsc{FlashAttention} proposed strategies for expediting attention on GPUs by optimizing for reduced memory access. Nonetheless, these techniques have not yet harnessed advanced features available in the latest hardware, with \textsc{FlashAttention-2} utilizing merely 35\% of the H100 GPU's computational capability. In this work, we introduce three pivotal advancements that accelerate attention on Hopper GPUs: (1) leveraging the asynchronous operations of both Tensor Cores and TMA to simultaneously execute computation and data transfer through warp-specialization, (2) fusing block-level matrix multiplication with softmax computations in an interleaved manner, and (3) employing block quantization alongside incoherent processing strategies enabled by hardware-accelerated FP8 precision. Our proposed approach, \textsc{FlashAttention-3}, boosts H100 GPU performance by factors of 1.5-2.0$\times$, attaining up to 740 TFLOPs/s (representing 75\% hardware utilization) with FP16 precision and approaching 1.2 PFLOPs/s when utilizing FP8. Moreover, we show that FP8 \textsc{FlashAttention-3} reduces numerical error by 2.6$\times$ in comparison to a baseline FP8 attention method.
\end{abstract}

\section{Introduction}
\label{sec:intro}

Within the Transformer architecture~\citep{vaswani2017attention}, the attention mechanism represents the main computational hurdle, primarily because calculating self-attention between queries and keys incurs a cost that grows quadratically with the sequence length. Extending attention to handle longer contexts could enable significant advances—such as processing and reasoning across multiple extended documents~\citep{guo2021longt5,shaham2022scrolls,peng2023yarn}, managing extensive codebases~\citep{roziere2023code, li2023starcoder}, handling new data modalities like high-resolution images~\citep{chen2022scaling}, audio~\citep{gulati2020conformer}, and video~\citep{ho2022video}, as well as supporting applications involving extended user histories~\citep{sun2019bert4rec} or prolonged agent workflows~\citep{yao2022react}. This challenge has spurred considerable research interest in improving the efficiency of attention operations for longer contexts, whether through approximation methods~\citep{katharopoulos2020transformers,choromanski2020rethinking,
tay2020efficient}, enhancements in software implementations~\citep{rabe2021self,dao2022flashattention,kwon2023efficient}, or by exploring fundamentally different model architectures~\citep{peng2023rwkv,sun2023retentive,gu2023mamba}.

This paper extends the approach of \citet{dao2022flashattention}, who devised exact-attention algorithms by thoughtfully incorporating the GPU’s architectural and execution traits into their algorithmic framework. In their prior work, \citet{dao2022flashattention} presented \textsc{FlashAttention}, introducing a novel tiling mechanism to parallelize attention in a manner that avoids unnecessary accesses to global memory; they accomplished this by combining all attention computations within a unified GPU kernel.

The later adaptation by \citet{dao2023flashattention2}, referred to as \textsc{FlashAttention-2}, modifies this scheme to leverage parallelism across the sequence length, restructuring the forward pass’s inner loops to operate on blocks of key and value matrices. This revision improved the spread and utilization of GPU resources. Despite these enhancements, our analysis reveals that \textsc{FlashAttention-2} still underperforms on recent GPU hardware when compared to high-efficiency matrix multiplication (GEMM) kernels — specifically, attaining only 35\% utilization versus the 80–90\% achievable on the Hopper H100 GPU. This disparity arises in part from differences at the implementation level, including the use of Ampere-era instructions rather than Hopper-optimized ones for operations on Tensor Cores. Recent contributions such as ThunkerKitten~\citep{spector2024thunder} and cuDNN 9~\citep{cudnn9} demonstrate that by employing Hopper-tailored instructions and employing tiling abstractions, attention mechanisms can be both accelerated and simplified in their codebase.

At a deeper level, the algorithm underlying \textsc{FlashAttention-2} operates within a straightforward synchronous paradigm, lacking any intentional integration of asynchronous execution or low-precision computations in its approach. Asynchrony emerges due to hardware architectures that target the acceleration of key components in machine learning workloads—dedicated units such as Tensor Cores focused on matrix multiplication, or the Tensor Memory Accelerator (TMA) responsible for efficient memory transfer, working independently of standard CUDA cores which handle general computation tasks involving logical operations and various numerical formats. The adoption of reduced-precision representations, including FP8 in Hopper and FP4 in Blackwell, with predecessors like FP16 (introduced with Pascal in 2017) and BF16 (featured in Ampere in 2020), has consistently enabled gains of twofold or greater throughput for equivalent power consumption and silicon real estate. The capabilities of Hopper in facilitating both asynchronous processing and low-precision are explored in detail in \cref{subsec:hardware}.

Redesigning \textsc{FlashAttention-2} to utilize these advanced hardware features presents several engineering hurdles: leveraging asynchrony necessitates strategies for overlapping the matrix multiplication and softmax computations, even when interdependencies exist, while incorporating low-precision arithmetic demands careful management of quantization artifacts, particularly concerning the problematic outlier activations encountered in large language models~\citep{dettmers2208llm, sun2024massive}.

In pursuit of this objective, we introduce \textsc{FlashAttention-3}, integrating and advancing three novel strategies aimed at optimizing performance on contemporary GPU architectures.\footnote{Our discussion primarily references NVIDIA's Hopper architecture, though the algorithm remains applicable to any GPU platform featuring strong asynchronous execution and support for low-precision computation.}

\begin{enumerate}

\item \textbf{Producer-Consumer asynchrony:} We introduce a warp-level software pipelining approach that takes advantage of the decoupled execution of data transfer and Tensor Core computations by separating the warps handling data production and consumption. This structural change enhances the pipeline's capacity to conceal both memory access delays and instruction issue latencies within the algorithm.

\item \textbf{Hiding softmax under asynchronous block-wise GEMMs:} We achieve overlap between softmax-related operations—such as floating point multiply-add and exponential—that are typically bandwidth-limited and the asynchronous WGMMA-based GEMM computations. To enable this, we redesign the \textsc{FlashAttention-2} algorithm, removing sequential dependencies that usually constrain execution order between softmax and GEMMs. As a practical instance, our 2-stage variant permits softmax processing on one block of the scores matrix concurrently with WGMMA computation of the following block, utilizing an asynchronous proxy.

\item \textbf{Hardware-accelerated low-precision GEMM:} Our forward pass procedure is adapted for use on FP8 Tensor Cores for GEMM, achieving nearly a twofold increase in empirical TFLOPs/s. This adaptation demands addressing distinct memory layout requirements imposed by WGMMA for FP32 accumulators and FP8 input matrices. To manage the accuracy degradation inherent to FP8 precision, we employ strategies such as block quantization and incoherent processing.
\end{enumerate}

To empirically substantiate our approach, we evaluate the performance of \textsc{FlashAttention-3} on the H100 SXM5 GPU across various configuration settings. The results indicate that (1) FP16 delivers a forward-pass speedup ranging from 1.5$\times$ to 2.0$\times$ over \textsc{FlashAttention-2} (achieving up to 740 TFLOPs/s), and backward-pass acceleration of 1.5-1.75$\times$, (2) FP8 attains nearly 1.2 PFLOPs/s throughput, and (3) for extensive sequence lengths, FP16 surpasses other methods, while FP8 demonstrates competitive performance\footnote{Specifically, with a head dimension of 64, FP8 in \textsc{FlashAttention-3} shows superior results; for head dimensions 128 and 256, FP8 is comparable to cuDNN attention when there is no causal masking but lags behind when causal masking is applied.} in comparison to NVIDIA’s cuDNN library's top-performing attention implementations. Furthermore, our findings confirm that FP16 in \textsc{FlashAttention-3} produces equivalent numerical errors as \textsc{FlashAttention-2} and exhibits improved accuracy relative to standard attention routines, since key intermediates such as softmax rescaling remain in FP32 precision. In addition, FP8 \textsc{FlashAttention-3} employing block quantization alongside incoherent processing demonstrates 2.6$\times$ higher accuracy than conventional attention employing per-tensor quantization, particularly in situations featuring outlier elements.

\textsc{FlashAttention-3} is released under a permissive license\footnote{\textsc{FlashAttention-3} can be found at \texttt{https://github.com/Dao-AILab/flash-attention}}, and we aim to incorporate its functionality into PyTorch and Hugging Face libraries so that it can be widely utilized by researchers and developers.

\section{Background: Multi-Head Attention and GPU Characteristics}
\label{sec:background}

\subsection{Multi-Head Attention}
\label{subsec:multi_head_attn}

Let $\mathbf{Q}, \mathbf{K}, \mathbf{V} \in \mathbb{R}^{N \times d}$ be the query, key and value input sequences associated to a single head, where $N$ is the sequence length and $d$ is the head dimension. Then the attention output $\mathbf{O}$ is computed as:

\begin{equation*}
  \mathbf{S} = \alpha \mathbf{Q} \mathbf{K}^\top \in \mathbb{R}^{N \times N}, \quad \mathbf{P} = \mathrm{softmax}(\mathbf{S}) \in \mathbb{R}^{N \times N}, \quad \mathbf{O} = \mathbf{P}\mathbf{V} \in \mathbb{R}^{N \times d},
\end{equation*}

Here, the $\mathrm{softmax}$ function operates on each row individually, and the scaling coefficient $\alpha$ is conventionally chosen as $1/\sqrt{d}$. To mitigate numerical issues associated with the exponential, it is standard to subtract $\mathrm{rowmax}(\mathbf{S})$ from $\mathbf{S}$. With respect to multi-head attention (MHA), distinct query, key, and value projection parameters are allocated for each attention head. This process is efficiently parallelized across all heads and batches to yield the composite output tensor.

Now let $\phi$ be a scalar loss function and let $\mathbf{d}(-) = \partial \phi / \partial (-)$ be notation for the gradient. Given the output gradient $\mathbf{dO} \in \mathbb{R}^{N \times d}$, we compute $\mathbf{dQ}$, $\mathbf{dK}$, and $\mathbf{dV}$ according to the chain rule as follows:

\begin{align*}
  \mathbf{dV} &= \mathbf{P}^\top \mathbf{dO} \in \mathbb{R}^{N \times d} \\
  \mathbf{dP} &= \mathbf{dO} \mathbf{V}^\top \in \mathbb{R}^{N \times N} \\
  \mathbf{dS} &= \mathrm{dsoftmax} (\mathbf{dP}) \in \mathbb{R}^{N \times N} \\
  \mathbf{dQ} &= \alpha \mathbf{dS} \mathbf{K} \in \mathbb{R}^{N \times d} \\
  \mathbf{dK} &= \alpha \mathbf{dS}^\top \mathbf{Q} \in \mathbb{R}^{N \times d},
\end{align*}

In this context, we observe that $\mathbf{d}s = (\mathrm{diag}(p) - p p^\top)\mathbf{d}p$, where $p$ is defined by the softmax function applied to the vector $s$. The notation $\mathrm{dsoftmax}(\mathbf{dP})$ denotes the row-wise application of this relationship. Moreover, this process is amenable to parallel execution over both attention heads and batches when performing the backward computation in MHA.

\subsection{GPU hardware characteristics and execution model}
\label{subsec:hardware}

We outline the components of the GPU execution model pertinent to \textsc{FlashAttention-3}, specifically emphasizing the NVIDIA Hopper architecture to exemplify these features.

\paragraph{Memory hierarchy:} The GPU architecture comprises multiple layers of memory, structured according to data location, where increased bandwidth is generally associated with reduced storage capacity (\cref{tab:gpu-hierarchy})\footnote{\citet{luo2024benchmarking} specifies a shared memory bandwidth of 128 bytes per clock cycle per SM; we compute aggregate bandwidth by scaling to 132 SMs and a boost frequency of 1830 MHz.}. At the top level, global memory (GMEM), or HBM, serves as the off-chip DRAM shared by all streaming multiprocessors (SMs). Data transfers from GMEM are automatically buffered by an on-chip L2 cache. Moving closer to the computation units, each SM has access to shared memory (SMEM)—a modestly sized, on-chip cache, whose access and allocation are managed explicitly by programmers. At the lowest level, SMs make use of their own register files.

\paragraph{Thread hierarchy:} In the GPU programming paradigm, threads are systematically arranged into several layers of logical groupings. These layers, ordered from smallest to largest, include individual threads, warps (each consisting of 32 threads), warpgroups (formed by 4 consecutive warps), threadblocks (also referred to as cooperative thread arrays or CTAs), threadblock clusters (introduced in Hopper architecture), and grids.

The interdependence of these two hierarchies is significant. Threads belonging to a single CTA are assigned concurrently to the same SM, while CTAs within a particular cluster are simultaneously dispatched to the same GPC. All threads in a CTA have shared access to SMEM, but each thread is limited to a maximum of 256 private registers (RMEM).

\paragraph{Asynchrony and warp-specialization:}

GPUs are designed as throughput-oriented processors that utilize parallelism and asynchronous operations to mask delays associated with memory access and instruction execution. To support asynchronous memory transfers between GMEM and SMEM, Hopper introduces the Tensor Memory Accelerator (TMA), a specialized hardware component \cite[\S7.29]{cuda}. In addition, the Hopper architecture distinguishes itself from previous designs like Ampere by providing an asynchronous Tensor Core accessible through the warpgroup-scoped WGMMA instruction \cite[\S9.7.14]{ptx}; notably, this Tensor Core is capable of ingesting data directly from shared memory.

The incorporation of hardware mechanisms for asynchrony enables the use of warp-specialized kernels, in which the warps within a CTA are designated either as producers or consumers, responsible exclusively for data transfer or computation, respectively. This division generally enhances the compiler's capability to produce more efficient instruction schedules \citep{warp-specialization-2011}. Furthermore, Hopper facilitates the dynamic redistribution of registers across warpgroups using \verb|setmaxnreg| \cite[\S9.7.17.1]{ptx}, thereby allowing warps engaged in MMAs to access a greater proportion of RMEM compared to warps that perform TMA operations, which require only a single thread.

\paragraph{Low-precision number formats:}
\label{sec:low-precision-gpu}
Contemporary GPUs feature dedicated hardware for efficient processing of reduced-precision calculations. As an illustration, the WGMMA instruction leverages the FP8 Tensor Cores present in Hopper GPUs, enabling twice the computational throughput per SM relative to FP16 or BF16 formats.

Nonetheless, proper use of FP8 WGMMA requires familiarity with the operand layout limitations. For a GEMM operation that computes $A \times B^{\top}$ with $A$ as an $M\times K$ matrix and $B$ as an $N\times K$ matrix, we refer to the $A$ or $B$ operand as \emph{mn-major} if the data is stored contiguously across the outer $M$ or $N$ dimension, respectively, and as \emph{k-major} if contiguity exists along the inner $K$ dimension instead. In the case of FP16 WGMMA, either mn-major or k-major formats are accepted for input operands residing in SMEM. By contrast, FP8 WGMMA only accommodates operands in the k-major format. Additionally, certain workloads such as attention mechanisms that require chaining multiple GEMM operations within a single kernel may run into difficulty: mismatched layouts between FP32 accumulators and FP8 operands prevent seamless use of subsequent dependent FP8 WGMMAs.

Within the framework of attention mechanisms, these constraints on layout necessitate particular adjustments to how an FP8 algorithm is constructed; further details are provided in \cref{sec:algofp8}.

\subsection{Standard Attention and Flash Attention} As defined by \citet{dao2022flashattention}, \textbf{standard attention} refers to a GPU-based attention implementation that explicitly stores the intermediate matrices $\mathbf{S}$ and $\mathbf{P}$ in HBM. The principal contribution of \textsc{FlashAttention} was to eliminate the overhead associated with these intermediate memory accesses by utilizing a localized form of the softmax reduction, thereby enabling attention computation within a single fused kernel. This local softmax operation is carried out in lines \ref{code-ws:softmax_start}-\ref{code-ws:softmax_end} of the consumer mainloop in \cref{alg:flash3_wgmma_ws_only}, accompanied by the appropriate rescaling of blocks within $\mathbf{O}$. A concise proof that this approach correctly produces $\mathbf{O}$ is presented in \cite[\S 2.3.1]{dao2023flashattention2}.

\section{FlashAttention-3: Algorithm}
\label{sec:algo}

This section presents the details of the \textsc{FlashAttention-3} algorithm. To maintain clarity, our focus remains on the forward pass, while the backward pass methodology is addressed separately in~\cref{sec:algo_ws_bwd}. We begin by illustrating the incorporation of warp-specialization alongside a circular SMEM buffer into the foundational framework of \textsc{FlashAttention-2}. Next, we discuss leveraging the asynchronous nature of WGMMA to build a GEMM-softmax pipeline featuring two overlapping stages. Lastly, we outline the necessary adaptations for FP8, encompassing both layout alignment and accuracy improvements through block quantization as well as incoherent processing.

\subsection{Producer-Consumer asynchrony through warp-specialization and pingpong scheduling}
\label{sec:algo_ws}

\paragraph{Warp-specialization}
Similar to \textsc{FlashAttention-2}, the forward computation in \textsc{FlashAttention-3} allows for high levels of parallelization across the batch dimension, number of attention heads, and the length of the query sequence. Accordingly, we focus on presenting the algorithm at the CTA (Cooperative Thread Array) level, where each CTA processes a particular tile $\mathbf{Q}_i$ from the query matrix and produces the associated tile $\mathbf{O}_i$ in the output. For clarity, our initial discussion centers on a warp-specialization strategy that utilizes a circular SMEM buffer, intentionally omitting any GEMM-softmax overlap at this stage. Denote the head dimension by $d$ and the sequence length by $N$. Choose a block size $B_r$ for partitioning $\mathbf{Q}$, which results in $T_r = \lceil \frac{N}{B_r} \rceil$ blocks, labeled as $\mathbf{Q}_1, .., \mathbf{Q}_{T_r}$.

\begin{algorithm}[H]
    \caption{\small\label{alg:flash3_wgmma_ws_only}\textsc{FlashAttention-3} forward pass \textbf{without} intra-consumer overlapping -- CTA perspective}
    \begin{algorithmic}[1]
\REQUIRE Matrices $\mathbf{Q}_i \in \mathbb{R}^{B_r \times d}$ and $\mathbf{K}, \mathbf{V} \in \mathbb{R}^{N \times d}$ residing in HBM, key block size $B_c$, and $T_c = \lceil \frac{N}{B_c} \rceil$.
\STATE Instantiate the pipeline manager responsible for barrier synchronization, utilizing an $s$-stage circular SMEM buffering scheme.
\IF {assigned to the producer warpgroup}
\STATE Release the predetermined registers as specified.
\STATE Begin transferring $\mathbf{Q}_i$ from HBM to shared memory.
\STATE After the load finishes, notify consumers that $\mathbf{Q}_i$ is now accessible.
\FOR{$0 \le j < T_c$}
    \STATE Await consumption of the $(j\,\%\,s)$ buffer stage by consumers.
    \STATE Load $\mathbf{K}_j, \mathbf{V}_j$ from HBM into the shared memory segment corresponding to the $(j\,\%\,s)$ buffer stage.
    \STATE Once complete, signal that $\mathbf{K}_j$ and $\mathbf{V}_j$ are available for consumers.
\ENDFOR
\ELSE \STATE Reserve a set number of registers for reuse, determined by the quantity of consumer warps present.
\STATE On-chip, set initial values: $\mathbf{O}_i = (0) \in \mathbb{R}^{B_r \times d}$, $\ell_i, m_i = (0), (-\infty) \in \mathbb{R}^{B_r}$.
\STATE Wait for $\mathbf{Q}_i$ to arrive in shared memory.
\FOR{$0 \le j < T_c$}
\STATE Wait until $\mathbf{K}_j$ has been moved to shared memory.
\STATE Carry out matrix multiplication $\mathbf{S}_i^{(j)} = \mathbf{Q}_i \mathbf{K}_j^T$ (SS-GEMM). Synchronize before proceeding. \label{alg:ws_only_gemm1}
\STATE Save $m_i^{\mathrm{old}} = m_i$ and update $m_i = \mathrm{max}(m_i^{\mathrm{old}}, \mathrm{rowmax}(\mathbf{S}_i^{(j)}))$. \label{code-ws:softmax_start}
\STATE Determine $\widetilde{\mathbf{P}}_i^{(j)} = \mathrm{exp}(\mathbf{S}_i^{(j)} - m_i)$ and update $\ell_i = \mathrm{exp}(m_i^{\mathrm{old}} - m_i) \ell_i + \mathrm{rowsum}(\widetilde{\mathbf{P}}_i^{(j)})$. \label{code-ws:softmax_end}
\STATE Wait until $\mathbf{V}_j$ is accessible in shared memory.
\STATE Compute $\mathbf{O}_i = \mathrm{diag}(\mathrm{exp}(m_i^{\mathrm{old}} - m_i))^{-1} \mathbf{O}_i + \widetilde{\mathbf{P}}_i^{(j)} \mathbf{V}_j$ (RS-GEMM). Synchronize and wait. \label{alg:ws_only_gemm2}
\STATE Free the $(j\,\%\,s)$ buffer segment to allow producer access.
\ENDFOR
\STATE Normalize $\mathbf{O}_i = \mathrm{diag}(\ell_i)^{-1} \mathbf{O}_i$ and compute $L_i = m_i + \log(\ell_i)$.
\STATE Store both $\mathbf{O}_i$ and $L_i$ in HBM, assigning them to the $i$th block in $\mathbf{O}$ and $L$ respectively.
\ENDIF
\end{algorithmic}
\end{algorithm}

In our deployment of \cref{alg:flash3_wgmma_ws_only} on Hopper, we utilize \verb|setmaxnreg| for memory (de)allocation purposes, employ TMA for loading $\mathbf{Q}_i$ and $\{ \mathbf{K}_j, \mathbf{V}_j \}_{0 \leq j < T_c}$, and rely on WGMMA for performing the GEMMs within the consumer mainloop. The SS or RS prefix designates whether the first GEMM operand originates from shared memory or the register file. To better understand the execution sequence in \cref{alg:flash3_wgmma_ws_only}, it is important to recognize that TMA load instructions operate asynchronously, meaning they are not held up by the completion of prior load operations. Additionally, within the producer mainloop, the initial $s$ iterations do not require any waits, as the buffer is still in the process of being populated.

\paragraph{Pingpong scheduling} Leveraging the asynchronous execution capabilities of WGMMA and TMA, together with warp specialization, enables the possibility of concurrently performing softmax calculations in one warpgroup while another warpgroup executes GEMM. This overlap is particularly desirable due to the disparity in throughput between matmul and non-matmul operations on current hardware accelerators. For instance, the H100 SXM5 GPU delivers 989 TFLOPS for FP16 matmul tasks but only achieves 3.9 TFLOPS for special functions such as exponential\footnote{According to the CUDA programming guide, each streaming multiprocessor (SM) supports 16 special function operations per clock cycle. With 132 SMs and a clock frequency of 1830 MHz, this equates to a total of 3.9 TFLOPS for special function operations.}, which are essential for softmax. During the FP16 attention forward pass with a head size of 128, the number of FLOPS required for matmul is 512 times greater than for exponentials; however, exponentials are computed at only 1/256 the throughput of matmuls, resulting in exponentials consuming up to 50\% of the execution cycles that matmuls use. This performance gap widens under FP8, where matmul throughput doubles but the exponential throughput remains unchanged.

The exponential operation is handled by a distinct hardware component—the multi-function unit—so, for optimal efficiency, it is preferable to overlap exponential computation with the matrix multiplication tasks executed by the Tensor Cores. To achieve this, we employ synchronization barriers (\texttt{bar.sync} instructions) to orchestrate the scheduling: the GEMMs (specifically, GEMM1, which computes $\mathbf{P} \mathbf{V}$ in the current iteration, and GEMM0, which computes $\mathbf{Q} \mathbf{K}^\top$ for the subsequent iteration) of warpgroup 1 are scheduled prior to those in warpgroup 2. Consequently, while warpgroup 2 is busy executing its GEMMs, the softmax operation associated with warpgroup 1 is underway. This scheduling pattern then reverses—warpgroup 2 performs the softmax while warpgroup 1 proceeds with its GEMMs, thus establishing a “pingpong” execution model. The mechanics of this approach are shown in~\cref{fig:pingpong_scheduling}. Although the real-world behavior of pingpong scheduling deviates somewhat from the idealized version presented in the diagram, empirical results indicate that this approach generally enhances throughput (for example, boosting performance from 570 TFLOPS to approximately 620–640 TFLOPS during FP16 forward passes with a head dimension of 128 and sequence length of 8192).

\paragraph{Attention variants} In the case of multi-query attention~\citep{shazeer2019fast} and grouped query attention~\citep{ainslie2023gqa}, we adopt the strategy utilized in \textsc{FlashAttention-2}. Specifically, tensor indexing is modified so as to prevent repeated storage of $\mathbf{K}$ and $\mathbf{V}$ within the HBM.

\subsection{Intra-warpgroup overlapping GEMMs and softmax}

It is possible to achieve partial instruction overlap between the softmax computations and the GEMM operations within a single warpgroup. Below, we outline a method for accomplishing this.

In the attention algorithm, operations within the inner loop (main loop) have sequential dependencies that impede parallelization within a single iteration. For example, (local) softmax (lines \ref{code-ws:softmax_start} to \ref{code-ws:softmax_end}) relies on the output $\mathbf{S}_i^{(j)}$ of the first GEMM, while the second GEMM takes its result $\widetilde{\mathbf{P}}_i^{(j)}$ as an operand. Indeed, the wait statements in lines \ref{alg:ws_only_gemm1} and \ref{alg:ws_only_gemm2} of \cref{alg:flash3_wgmma_ws_only} serialize the execution of softmax and GEMMs. However, we can break these dependencies by pipelining across iterations through additional buffers in registers. Pursuing this idea, we propose the following two-stage\footnote{Note that the number of stages of the overlapping scheme is bounded by, but need not equal, the number $s$ of stages in the circular SMEM buffer.} GEMM-softmax pipelining algorithm:

\begin{algorithm}[H]
    \caption{\small\label{alg:flash3_wgmma}\textsc{FlashAttention-3} consumer warpgroup forward pass}
    \begin{algorithmic}[1]
      \REQUIRE Matrices $\mathbf{Q}_i \in \mathbb{R}^{B_r \times d}$ and $\mathbf{K}, \mathbf{V} \in \mathbb{R}^{N \times d}$ stored in HBM; key block size $B_c$, and $T_c = \lceil \frac{N}{B_c} \rceil$.
      \STATE Dynamically allocate registers in accordance with the count of consumer warps.
      \STATE Initialize in on-chip memory: $\mathbf{O}_i = (0) \in \mathbb{R}^{B_r \times d}$ as well as $\ell_i, m_i = (0), (-\infty) \in \mathbb{R}^{B_r}$.
      \STATE Synchronize until both $\mathbf{Q}_i$ and $\mathbf{K}_0$ are resident in shared memory.
      \STATE Calculate $\mathbf{S}_{\mathrm{cur}} = \mathbf{Q}_i \mathbf{K}_0^T$ using WGMMA; commit computation and await completion.
      \STATE Free the $0$th buffer stage allocated to $\mathbf{K}$.
      \STATE Determine values $m_i$, $\tilde{\mathbf{P}}_{\mathrm{cur}}$, and $\ell_i$ from $\mathbf{S}_{\mathrm{cur}}$; update scaling of $\mathbf{O}_i$.
      \FOR{$1 \le j < T_c - 1$}
        \STATE Hold until $\mathbf{K}_j$ is available in shared memory.
        \label{code:mainloop_start}
        \STATE Use WGMMA to perform $\mathbf{S}_{\mathrm{next}} = \mathbf{Q}_i \mathbf{K}_j^T$; commit task without waiting.
        \label{code:first_wgmma}
        \STATE Wait for $\mathbf{V}_{j-1}$ to arrive in shared memory.
        \STATE Execute $\mathbf{O}_i = \mathbf{O}_i + \tilde{\mathbf{P}}_{\mathrm{cur}} \mathbf{V}_{j-1}$ via WGMMA; commit without waiting.
        \label{code:second_wgmma}
        \STATE Await the completion of WGMMA computation $\mathbf{Q}_i \mathbf{K}_j^T$.
        \STATE Compute $m_{i}$, $\tilde{\mathbf{P}}_{\mathrm{next}}$, and $\ell_{i}$ from $\mathbf{S}_{\mathrm{next}}$.
        \label{code:softmax}
        \STATE Wait for WGMMA result $\tilde{\mathbf{P}}_{\mathrm{cur}} \mathbf{V}_{j-1}$ and then rescale $\mathbf{O}_i$.
        \STATE Release the stage $(j\,\%\,s)$ and $(j-1\,\%\,s)$ in the buffer corresponding to $\mathbf{K}$ and $\mathbf{V}$, respectively.
        \STATE Assign $\mathbf{S}_{\mathrm{cur}}$ the value of $\mathbf{S}_{\mathrm{next}}$.
        \label{code:mainloop_end}
      \ENDFOR
      \STATE Wait for $\mathbf{V}_{T_c - 1}$ to be present in shared memory.
      \STATE Calculate $\mathbf{O}_i = \mathbf{O}_i + \tilde{\mathbf{P}}_{\mathrm{last}} \mathbf{V}_{T_c - 1}$ employing WGMMA; commit and wait.

\STATE Finalization: Adjust $\mathbf{O}_{i}$ according to $m_{i}$. Determine $L_{i}$ from both $m_{i}$ and $\ell_{i}$.
Save $\mathbf{O}_{i}$ and $L_{i}$ in HBM as the $i$-th segments of $\mathbf{O}$ and $L$.

\cref{alg:flash3_wgmma} substitutes for the consumer pathway found in \cref{alg:flash3_wgmma_ws_only}, thereby constituting the full \textsc{FlashAttention-3} procedure tailored for FP16 computations. Broadly, the term WGMMA serves to denote asynchronous GEMM operations. During the mainloop, which spans lines \ref{code:mainloop_start} through \ref{code:mainloop_end}, the second WGMMA execution occurring in iteration $j$ (line \ref{code:second_wgmma}) is executed concurrently with the softmax computations corresponding to iteration $j+1$ (line \ref{code:softmax}).

Although the pipelined model depicted earlier presents potential performance advantages, several implementation details must be acknowledged:

\paragraph{Compiler reordering} The pseudocode reflects an ideal, sequential workflow; however, the NVCC compiler frequently reorganizes instructions to achieve better optimization. This restructuring can interfere with the intended ordering of WGMMA and non-WGMMA pipelined operations, possibly resulting in unpredictable outcomes or reducing the expected speedup. Examination of the SASS output demonstrates that the compiler does produce the intended overlapped instruction streams (see Section~\ref{sec:2-stage-sass}).

\paragraph{Register pressure} Achieving high performance demands careful management to avoid register spilling. The 2-stage pipeline architecture introduces extra register usage: it is necessary to reserve space for intermediate storage and continuity between pipeline stages. Specifically, storing an additional $\mathbf{S}_{\mathrm{next}}$ in registers is required, increasing register consumption by $B_r \times B_c \times \text{sizeof}(\text{float})$ per threadblock. This greater register need can limit feasible block sizes, which are themselves common optimization targets requiring substantial registers. As a result, these competing factors should be evaluated through empirical performance measurements and appropriate balancing.

\paragraph{3-stage pipelining} Building upon the preceding 2-stage scheme, we introduce a 3-stage method that further overlaps a second WGMMA operation with the softmax computation. This extension has the potential to improve Tensor Core throughput, yet it significantly escalates the register requirements, as one more pipeline stage must be maintained. Consequently, balancing pipeline depth and tile dimensions becomes increasingly complex. For more information about the algorithm and its performance assessment, refer to ~\cref{sec:3-stage}.

\subsection{Low-precision with FP8}
\label{sec:algofp8}

\textbf{Efficiency: layout transformations.} Executing the forward pass of \textsc{FlashAttention-3} using FP8 precision introduces extra difficulties related to layout compatibility, which are not present when working with FP16.

Initially, we observe that input tensors $\mathbf{Q}$, $\mathbf{K}$, and $\mathbf{V}$ are conventionally stored with the head dimension as the contiguous axis. However, the FP8 WGMMA's k-major requirement for the second GEMM necessitates that $\mathbf{V}$—specifically, the tiles loaded into SMEM—be contiguous along the sequence length dimension. Given that TMA loads cannot alter which dimension is contiguous, this presents two alternatives: (1) perform a transpose of $\mathbf{V}$ in GMEM beforehand, or (2) execute an in-kernel transpose on the $\mathbf{V}$ tiles after their transfer to SMEM. Option (1) could be approached by either (1a) fusing the transpose operation with the epilogue of a previous stage, such as the rotary embedding, or (1b) deploying a dedicated pre-processing transpose kernel\footnote{Optimized transpose kernels can attain performance approaching device memory bandwidth~\citep{colfax_cutlass_transpose_2024}.} to swap the strides between sequence length and head dimensions. Yet, integrating (1a) with standard library workflows proves challenging, and (1b) can lead to excessive overhead in memory-constrained environments typical of inference workloads.

For FP8 \textsc{FlashAttention-3}, we choose to utilize approach (2). The in-kernel transposition is facilitated by making use of LDSM (\verb|ldmatrix|) and STSM (\verb|stmatrix|) instructions, which enable a warp of threads to collaboratively move data between SMEM and RMEM, operating on 128-byte blocks.\footnote{According to the PTX documentation, LDSM/STSM are intended to transfer $8 \times 8$ matrices containing 16-bit elements \cite[\S 9.7.13.4.15-16]{ptx}. Nonetheless, by packing two 8-bit elements together, it becomes feasible to apply LDSM/STSM instructions within the framework of FP8 arithmetic. Yet, the transpose variants of LDSM/STSM cannot independently access packed 8-bit values, prompting the need for additional register shuffling between LDSM and STSM invocations in order to realize a transposed tile; we refrain from providing the explicit method here.} Not only are these instructions frugal with register resources, which permits execution within the producer warpgroup, but they also can effect layout transposition during the memory copy operation. Furthermore, after the initial iteration, it is possible to schedule the transposition of the subsequent $\mathbf{V}$ tile to occur concurrently with the two WGMMAs that process the previous $\mathbf{V}$ and current $\mathbf{K}$ tiles.

Secondly, we note that, in contrast to FP16, the FP32 accumulator's memory layout for an FP8 WGMMA does not correspond to the layout expected for operand A when stored in registers. Partial depictions of these layouts are presented in \cref{fig:rmem_accum} and \cref{fig:rmem_operand}, with register entries per thread arranged according to the order indicated. Utilizing byte permute instructions enables us to convert the accumulator resulting from the initial WGMMA into an arrangement compatible with the subsequent WGMMA, as well as aligning with the $\mathbf{V}$ tile layout generated by the in-kernel transpose. More precisely, referencing \cref{fig:rmem_accum}, we reorder the sequence to
$$\{ \verb|d0 d1 d4 d5 d2 d3 d6 d7| \},$$
and repeat this register permutation pattern every 8 bytes. When considering the logical structure of the $\mathbf{P}$ tile, this operation performs a column permutation—shifting, for instance, columns $0189$ to the leading column positions. To ensure WGMMA produces the correct output tile, the in-kernel transpose can be coordinated to output a corresponding row permutation of the $\mathbf{V}$ tile.\footnote{This ability to leverage the in-kernel transpose in this way obviates the necessity for shuffle instructions to reassign register ownership between threads, as discussed previously in~\citep{colfax_fp8_flashattention_2024}.}

\textbf{Accuracy: block quantization and incoherent processing.} The FP8 (e4m3) representation encodes each value with 3 mantissa bits and 4 exponent bits, leading to increased quantization noise compared to FP16 or BF16. Furthermore, outlier elements—those with significantly greater magnitude—are often present in large-scale models~\citep{dettmers2208llm, sun2024massive}, complicating the quantization process. Standard practice employs per-tensor scaling~\citep{micikevicius2022fp8}, assigning a single scaling factor to each tensor (e.g., separately for $\mathbf{Q}$, $\mathbf{K}$, and $\mathbf{V}$). To mitigate the accuracy degradation in attention caused by FP8 quantization, we introduce two methods:

\begin{enumerate}

\item \textbf{Block quantization}: this approach assigns a unique scalar to each block within the tensor, segmenting $\mathbf{Q}$, $\mathbf{K}$, and $\mathbf{V}$ into sub-tensors of dimensions $B_r \times d$ or $B_c \times d$ and independently quantizing each block. The quantization process can be seamlessly integrated with preceding operations such as rotary embedding without increasing runtime, due to the memory-bandwidth constraints of the embedding. Given that the \textsc{FlashAttention-3} algorithm intrinsically manages data in blocks, adjusting each block of $\mathbf{S}$ for quantization can be accomplished without extra computational overhead.

\item \textbf{Incoherent processing}: to mitigate the influence of anomalous values, $\mathbf{Q}$ and $\mathbf{K}$ are pre-multiplied by a random orthogonal matrix $\mathbf{M}$ before being quantized to FP8. The orthogonality condition $\mathbf{M} \mathbf{M}^\top = I$ ensures that the transformed attention calculation $(\mathbf{Q} \mathbf{M}) (\mathbf{K} \mathbf{M})^\top$ remains equivalent to the original $\mathbf{Q} \mathbf{K}^\top$. This randomization redistributes the outliers, as each element of $\mathbf{Q} \mathbf{M}$ or $\mathbf{K} \mathbf{M}$ represents a random combination of elements from the original matrices, thereby decreasing the error introduced by quantization. Following the methodology of \citet{chee2024quip} and~\citet{tseng2024quip}, $\mathbf{M}$ is constructed as a product of randomly signed diagonal matrices and a Hadamard matrix. This choice allows multiplication to be performed in $O(d \log d)$ time rather than $O(d^2)$, and enables fusion with rotary embedding with no extra computational requirements.

We demonstrate in \cref{sec:numerical_error} that applying these methods decreases numerical error by as much as 2.6$\times$.

\section{Empirical Validation}
\label{sec:exp}

To implement \textsc{FlashAttention-3} and assess its performance and correctness, we leverage the WGMMA and TMA abstractions provided by the CUTLASS~\citep{Thakkar_CUTLASS_2023} library.

\begin{itemize}

\item \textbf{Benchmarking attention.} We assess the performance of \textsc{FlashAttention-3} by recording its runtime over a range of sequence lengths, contrasting the results with those from both the baseline PyTorch implementation, \textsc{FlashAttention-2}, the Triton-based version of \textsc{FlashAttention-2} (leveraging H100-specific instructions), and a cuDNN-based vendor-optimized \textsc{FlashAttention-2} for H100 GPUs. Our findings reveal that \textsc{FlashAttention-3} achieves up to a 2.0$\times$ speedup over \textsc{FlashAttention-2}, as well as a 1.5$\times$ improvement compared to the Triton implementation. Additionally, \textsc{FlashAttention-3} attains peak throughput of 740 TFLOPs/s, corresponding to 75\% of the peak TFLOPs/s achievable on H100 GPUs.
\item \textbf{Ablation study.} We demonstrate that optimizations such as warp-specialization and the incorporation of GEMM-softmax pipelining are principal factors driving the increased efficiency observed in \textsc{FlashAttention-3}.
\item \textbf{Accuracy of FP8 attention.} Through experimentation, we show that utilizing block quantization alongside incoherent processing techniques effectively reduces the numerical error in the FP8 implementation of \textsc{FlashAttention-3} by a factor of 2.6$\times$.
\end{itemize}

\subsection{Benchmarking Attention}
\label{subsec:benchmark_attn}

We evaluate the execution time of various attention mechanisms using an H100 80GB SXM5 GPU across different conditions (including the presence or absence of a causal mask and head dimensions of 64 or 128) with FP16 input representations. The outcomes of these benchmarks are presented in~\cref{fig:benchmark_attn_fwd} and~\cref{fig:benchmark_attn_bwd}. The data indicate that \textsc{FlashAttention-3} achieves approximately 1.5--2.0$\times$ speedup over \textsc{FlashAttention-2} during forward computation and about 1.5--1.75$\times$ speedup during backward computation. Relative to conventional attention implementations, \textsc{FlashAttention-3} delivers improvements in runtime of up to 3--16$\times$. Notably, for sequence lengths of 1,000 tokens and above, \textsc{FlashAttention-3} even outperforms the vendor-provided cuDNN library—which is specifically optimized for the H100 GPU platform—despite cuDNN being closed source.

\paragraph{Benchmark settings:} We vary the sequence length as 512, 1k, ..., 16k, and set batch size so that the total number of tokens is 16k. We set the hidden dimension to 2048, and head dimension to be either 64, 128, or 256 (i.e., 32 heads, 16 heads, or 8 heads). To calculate the FLOPs of the forward pass, we use:

\begin{equation*}
  4 \cdot \text{seqlen}^2 \cdot \text{head dimension} \cdot \text{number of heads}.
\end{equation*}

By implementing causal masking, we halve the total value to reflect that roughly 50% of the positions are computed. To estimate the FLOPs for the backward pass, we scale the forward pass FLOPs by a factor of 2.5, as the backward operation involves five matrix multiplications, compared to two in the forward pass, resulting from recomputation.

Additionally, we evaluate the forward pass runtime utilizing FP8 under comparable conditions. The outcomes for headdim 256 are illustrated in ~\cref{fig:benchmark_attn_fp8}, while comprehensive results are provided in \cref{sec:benchmark_attn_fp8_full}.

\subsection{Ablation Study: 2-Stage Pipelining Experiments}
\label{sec:2-stage-pipelining-experiments}

We conduct an ablation study by disabling both the 2-stage WGMMA-softmax pipeline and warp-specialization features for non-causal FP16 \textsc{FlashAttention-3}, maintaining the configuration $\{\text{batch}, \text{seqlen}, \text{nheads}, \text{hdim}\} = \{ 4, 8448, 16, 128\}$. The findings shown in~\cref{table:ablation_pipelining} demonstrate that the introduced algorithmic strategies—namely, asynchronous execution via warp-specialization and the concurrent processing of GEMM and softmax—yield notable performance improvements, boosting throughput from 570 to 661 TFLOPs.

\subsection{Numerical Error Validation}
\label{sec:numerical_error}

Given the ongoing discussion regarding the numerical accuracy of \textsc{FlashAttention}~\citep{golden2024flash}, we evaluate and contrast \textsc{FlashAttention-2}, \textsc{FlashAttention-3}, as well as a conventional attention implementation, with a high-precision FP64 reference. To emulate the presence of outlier activations and features characteristic of LLMs~\citep{dettmers2208llm, sun2024massive}, we sample $\mathbf{Q}, \mathbf{K}, \mathbf{V}$ according to the following statistical distribution:

\begin{equation*}
  \mathcal{N}(0, 1) + \mathcal{N}(0, 100) \cdot \mathrm{Bernoulli}(0.001).
\end{equation*}

In this setup, every entry follows a normal distribution with mean zero and unit variance, except for 0.1\% of the entries, where we introduce an additional independent term drawn from a normal distribution with standard deviation 10. We report the root mean squared error (RMSE) in~\cref{table:numerical_error}. For FP16, both \textsc{FlashAttention-2} and \textsc{FlashAttention-3} demonstrate an RMSE that is 1.7$\times$ lower than that of the standard implementation, as they maintain intermediate softmax computations in FP32. The baseline attention method with FP8 applies per-tensor scaling and stores the matmul accumulator in FP32, while the softmax intermediates are preserved in FP16. Due to the use of block quantization and incoherent processing, \textsc{FlashAttention-3} in FP8 outperforms this baseline, achieving 2.6$\times$ the accuracy.

\section{Dicussion, Limitations, Conclusion}
\label{sec:discussion}

Through the development of \textsc{FlashAttention-3}, we show that innovative programming approaches and hardware capabilities—particularly asynchrony and reduced-precision computation—can substantially enhance the performance and reliability of attention mechanisms. Our methods achieve a 1.5-2.0$\times$ improvement in attention speed relative to \textsc{FlashAttention-2}, and decrease FP8 numerical error by 2.6$\times$ compared to conventional per-tensor quantization strategies. Future work will focus on several outstanding challenges: enhancing LLM inference, implementing a persistent kernel architecture within the FP8 kernel,\footnote{In our current evaluations, FP16 \textsc{FlashAttention-3} utilizes a persistent kernel and a load balancing scheme, while FP8 \textsc{FlashAttention-3} lacks these features. This difference partially accounts for the inferior performance of FP8 \textsc{FlashAttention-3} for short sequences and under causal masking, when contrasted with FP8 cuDNN kernels.} and investigating the implications of low-precision attention in large-scale training scenarios. Although our experiments are limited to Hopper GPUs, we anticipate that the underlying methods will extend to a variety of hardware accelerators. Ultimately, we hope that improvements in attention speed and fidelity will open up further possibilities for applications requiring long-context processing.

\end{document}